%Version 3.1 December 2024
% Springer Nature LaTeX Template (sn-jnl) customized for Nexus Project Report

\documentclass[pdflatex,sn-mathphys-num]{sn-jnl}% Math and Physical Sciences Numbered Reference Style

%%%% Standard Packages
\usepackage{graphicx}%
\usepackage{multirow}%
\usepackage{amsmath,amssymb,amsfonts}%
\usepackage{amsthm}%
\usepackage{mathrsfs}%
\usepackage[title]{appendix}%
\usepackage{xcolor}%
\usepackage{textcomp}%
\usepackage{manyfoot}%
\usepackage{booktabs}%
\usepackage{algorithm}%
\usepackage{algorithmicx}%
\usepackage{algpseudocode}%
\usepackage{listings}%
\usepackage{hyperref}%

% Code snippet styling
\definecolor{codegreen}{rgb}{0,0.6,0}
\definecolor{codegray}{rgb}{0.5,0.5,0.5}
\definecolor{codepurple}{rgb}{0.58,0,0.82}
\definecolor{backcolour}{rgb}{0.96,0.96,0.96}

\lstdefinestyle{mystyle}{
    backgroundcolor=\color{backcolour},   
    commentstyle=\color{codegreen},
    keywordstyle=\color{magenta},
    numberstyle=\tiny\color{codegray},
    stringstyle=\color{codepurple},
    basicstyle=\ttfamily\scriptsize,
    breakatwhitespace=false,         
    breaklines=true,                 
    captionpos=b,                    
    keepspaces=true,                 
    numbers=left,                    
    numbersep=5pt,                  
    showspaces=false,                
    showstringspaces=false,
    showtabs=false,                  
    tabsize=2
}
\lstset{style=mystyle}

%% Theorem styles
\theoremstyle{thmstyleone}%
\newtheorem{theorem}{Theorem}%
\newtheorem{proposition}[theorem]{Proposition}%

\theoremstyle{thmstyletwo}%
\newtheorem{example}{Example}%
\newtheorem{remark}{Remark}%

\theoremstyle{thmstylethree}%
\newtheorem{definition}{Definition}%

\raggedbottom

\begin{document}

\title[Nexus: AI-Powered Semantic Desktop Search Engine]{Nexus: An AI-Powered Semantic Desktop Search Engine Utilizing Hybrid Dense-Sparse Retrieval and a Decoupled Sidecar Architecture}

\author[1]{\fnm{Aadhisesha} \sur{D}}
\author[1]{\fnm{Roshan Kumar} \sur{K}}

\affil[1]{\orgdiv{Department of Computer Science and Engineering}, \orgname{College of Engineering, Guindy}, \orgaddress{\city{Chennai}, \postcode{600025}, \state{Tamil Nadu}, \country{India}}}

%%==================================%%
%% Abstract                         %%
%%==================================%%

\abstract{Traditional desktop file search utilities embedded within modern operating systems (such as macOS Spotlight and Windows Search) rely predominantly on exact lexical pattern matching and surface-level filesystem metadata. Consequently, they suffer from the ``vocabulary mismatch problem'', lack conceptual understanding, and exhibit near-total blindness to unstructured visual media. Conversely, cloud-based enterprise cognitive search engines introduce severe data sovereignty, bandwidth, and confidentiality risks when deployed across local personal workstations. In this paper, we present \textbf{Nexus}, an extensible, privacy-first desktop search engine engineered to provide high-precision semantic discovery across heterogeneous local document corpora. Nexus introduces a decoupled desktop architecture: a lightweight Tauri (Rust) shell manages the native presentation layer (React + Vite + TypeScript) and operating system security capabilities, while a supervised Python/FastAPI sidecar service executes asynchronous document ingestion, tokenization, and vector indexing. To balance precision and recall, Nexus pairs sparse lexical retrieval (Okapi BM25) with dense semantic vector representations stored in an embedded vector database (ChromaDB), harmonized through convex score combination and Reciprocal Rank Fusion (RRF). Furthermore, Nexus implements an abstract multi-provider embedding layer capable of orchestrating hosted foundation models (e.g., OpenAI \texttt{text-embedding-3}, Google Gemini, Anthropic) and multimodal vision-language models for automated image captioning. Crucially, Nexus operates under a zero-trust credential model wherein user-supplied API keys are encrypted at rest using authenticated AES-256-GCM and transmitted strictly to direct model provider endpoints without intermediary proxy telemetry. We articulate the system design, mathematical retrieval formulation, and an operational evaluation demonstrating substantial memory and latency advantages over conventional desktop frameworks.}

\keywords{Semantic Search, Hybrid Retrieval, Okapi BM25, Dense Embeddings, Tauri Sidecar, ChromaDB, Desktop Privacy, Multimodal Ingestion}

\maketitle

% =========================================================================
% SECTION 1: INTRODUCTION
% =========================================================================
\section{Introduction}\label{sec1}

The exponential growth of digital documents on personal and workstation computers has rendered effective local information discovery increasingly arduous. Users routinely maintain hundreds of gigabytes of heterogeneous files, including Portable Document Format (PDF) files, rich office documents (\texttt{.docx}, \texttt{.xlsx}, \texttt{.pptx}), source code repositories, and unstructured raster imagery. Despite major advances in consumer hardware, default operating system search systems---such as Windows Search, macOS Spotlight, and Linux \texttt{locate}/Recoll---continue to rely primarily on inverted lexical indices, exact token matching, and basic filesystem attributes (e.g., file modification timestamp, file size, and filename strings)~\cite{baeza1999modern,manning2008introduction}.

This paradigm introduces three fundamental systemic bottlenecks:
\begin{enumerate}
    \item \textbf{The Vocabulary Mismatch Problem}: Lexical matching cannot resolve synonyms, paraphrased intent, or conceptual abstraction. A query such as \textit{``quarterly financial liabilities''} will fail to retrieve a file titled \textit{``Q3 Debt Obligations.xlsx''} if exact term overlap is absent~\cite{furnas1987vocabulary}.
    \item \textbf{Visual and Multimodal Blindness}: Image files, architectural diagrams, scanned receipts, and interface mockups are indexed solely by arbitrary filename strings unless costly optical character recognition (OCR) is explicitly executed. Conceptual elements depicted within images remain opaque to text-based lexical parsers.
    \item \textbf{Privacy versus Capability Duality}: While cloud enterprise cognitive search engines (e.g., Microsoft Copilot for Enterprise, Google Cloud Vertex Search) solve semantic understanding through deep Transformer-based encoders~\cite{vaswani2017attention}, uploading entire confidential hard drives to third-party multi-tenant cloud data stores presents catastrophic data sovereignty, confidentiality, and regulatory non-compliance risks (e.g., under GDPR and HIPAA)~\cite{voigt2017eu}.
\end{enumerate}

To bridge this gap, this paper introduces \textbf{Nexus}, an open, privacy-centric, desktop-native semantic search engine. Unlike preceding monolithic client implementations, Nexus adopts a decoupled, multi-runtime architecture. The visual desktop environment is powered by Tauri, a Rust-based desktop framework with minimal memory overhead, while the heavy numerical and machine learning computations are offloaded to an asynchronous Python and FastAPI sidecar process.

\noindent\textbf{Key Contributions.} The primary contributions of this work are summarized as follows:
\begin{itemize}
    \item \textbf{Decoupled Sidecar Architecture}: We design and implement a dual-runtime desktop topology utilizing Tauri (Rust) as a secure desktop shell and a localized FastAPI service running in a background sidecar. This eliminates the heavy RAM footprint of conventional Electron wrappers while maintaining access to Python's rich data science and natural language processing ecosystems.
    \item \textbf{Hybrid Dense-Sparse Retrieval Engine}: We formulate an adaptive retrieval engine unifying BM25 sparse lexical scoring with dense semantic embedding distances stored in an embedded vector store (ChromaDB), mitigating the weaknesses of pure vector search on keyword-specific code identifiers.
    \item \textbf{Multimodal Ingestion Pipeline}: We formalize an extensible ingestion pipeline that incorporates multimodal vision-language captioning to convert raster imagery into rich semantic vectors.
    \item \textbf{Zero-Trust Credential and Privacy Model}: We deploy client-side authenticated key storage using symmetric AES-256-GCM encryption with machine-derived entropy, guaranteeing that raw documents and search tokens never traverse intermediary proxy services.
\end{itemize}

% =========================================================================
% SECTION 2: RELATED WORK
% =========================================================================
\section{Related Work and Theoretical Foundations}\label{sec2}

\subsection{Lexical versus Dense Semantic Retrieval}\label{subsec2_1}
Classical information retrieval is grounded in statistical term weighting. The Okapi BM25 algorithm~\cite{robertson2009probabilistic} computes non-linear term frequencies normalized against document length. While extraordinarily robust for exact identifier lookup (such as function names, error codes, and unique proper nouns), BM25 cannot capture semantic affinity across differing vocabularies.

The emergence of pretrained Transformer encoders~\cite{devlin2018bert} and dense bi-encoder models such as Sentence-BERT (SBERT)~\cite{reimers2019sentence} enabled semantic projection into continuous metric spaces where cosine similarity correlates directly with semantic relatedness:
\begin{equation}
\text{sim}(\mathbf{u}, \mathbf{v}) = \frac{\mathbf{u} \cdot \mathbf{v}}{\|\mathbf{u}\|_2 \|\mathbf{v}\|_2}\label{eq_cosine}
\end{equation}
However, pure dense retrieval frequently suffers from catastrophic forgetting regarding out-of-vocabulary technical identifiers, version numbers, and verbatim quotes. Recent literature emphasizes hybrid retrieval~\cite{karpukhin2020dense,lin2021pretrained} as the optimal consensus, harmonizing exact syntactic fidelity with latent conceptual understanding. Crucially, Zerhoudi et al.~\cite{zerhoudi2026aswemaysearch} formally established the foundations of ``local-first information retrieval'', demonstrating that on-device dense and hybrid pipelines operating over consumer hardware maintain upwards of 91\% nDCG@10 up to 100K personal documents with zero privacy leakage. Nexus builds directly upon these empirical findings by engineering a concrete desktop-native hybrid search deployment.

\subsection{Desktop AI Application Architectures}\label{subsec2_2}
Historically, desktop applications integrating modern web graphical interfaces have depended on the Chromium Embedded Framework (CEF) and Electron. However, Electron bundles complete Chromium and Node.js binaries, routinely consuming hundreds of megabytes of idle memory~\cite{tauri2023docs}. In response, modern system-native frameworks such as Tauri leverage native platform webviews (WebKit on macOS, WebView2 on Windows) bound to Rust binaries. By managing external processing engines through managed child processes (sidecars), applications achieve near-instantaneous startup times, rigorous permission sandboxing, and reduced execution overheads.

% =========================================================================
% SECTION 3: SYSTEM ARCHITECTURE
% =========================================================================
\section{System Architecture and Topology}\label{sec3}

Nexus is architected around the principle of physical runtime separation between interface rendering, operating system integration, and computational machine learning pipelines. Table~\ref{tab_arch} outlines the structural distribution of responsibilities across the application tiers.

\begin{table}[h]
\caption{Architectural decomposition and layer responsibilities in Nexus.}\label{tab_arch}
\begin{tabular*}{\textwidth}{@{\extracolsep\fill}lll@{}}
\toprule
\textbf{Layer} & \textbf{Core Technologies} & \textbf{Operational Responsibilities} \\
\midrule
Presentation Layer & React 18, Vite, TypeScript & Search input, virtualized result streams, settings modal \\
Desktop Host Shell & Tauri 2.0 (Rust Kernel) & Sidecar life-cycle supervision, native dialogs, sandboxing \\
IPC Bridge & Loopback HTTP / WebSocket & Local non-routable communication (\texttt{http://127.0.0.1:<port>}) \\
Computation Engine & Python 3.11, FastAPI & Async document parsing, token chunking, embedding dispatch \\
Storage Layer & ChromaDB, SQLite & HNSW vector indices, document metadata and chunk hashes \\
\botrule
\end{tabular*}
\end{table}

\subsection{Host Desktop Shell and Sidecar Orchestration}\label{subsec3_1}
The host layer is implemented in Rust utilizing the Tauri framework. Tauri serves three distinct operational roles:
\begin{itemize}
    \item \textbf{Process Supervisor}: Upon application launch, the Rust host spawns the compiled Python sidecar binary as a managed child subprocess. The host continuously monitors process liveness, reads standard error/output streams for telemetry, and guarantees graceful termination upon desktop window closure.
    \item \textbf{Operating System Capability Boundary}: Native filesystem directory selection dialogs and security permissions are enforced natively in Rust via declarative capabilities defined in \texttt{tauri.conf.json}.
    \item \textbf{Secure Webview Environment}: The user interface is isolated from arbitrary shell execution, preventing malicious document injections from executing arbitrary code.
\end{itemize}

\subsection{Presentation Layer}\label{subsec3_2}
The user interface is engineered in React 18 with Vite and TypeScript. It utilizes a declarative state model to manage:
\begin{enumerate}
    \item \textbf{Asynchronous Query Dispatch}: Queries entered by the user are debounced and transmitted via an internal API abstraction client to the local FastAPI daemon.
    \item \textbf{Streaming Retrieval Updates}: Results are rendered incrementally using virtualized lists to maintain a fluid 60 frames-per-second (FPS) rendering loop even when dealing with thousands of document chunks.
    \item \textbf{Configuration Management}: A localized settings dashboard manages provider configuration (e.g., OpenAI, Google Gemini, Anthropic) and credential provisioning.
\end{enumerate}

\subsection{FastAPI Sidecar Engine}\label{subsec3_3}
The core computation engine executes within a Python runtime encapsulated via FastAPI and Uvicorn. By employing Python's asynchronous event loop (\texttt{asyncio}), the sidecar performs non-blocking ingestion of large filesystems while concurrently servicing interactive user search requests. The sidecar exposes standardized REST endpoints:
\begin{itemize}
    \item \texttt{GET /health}: System status and vector store indexing progress.
    \item \texttt{POST /search}: Dispatches hybrid retrieval queries across indices.
    \item \texttt{POST /index}: Accepts filesystem directories for recursive crawling.
    \item \texttt{POST /settings}: Ingests and verifies encrypted provider credentials.
\end{itemize}

% =========================================================================
% SECTION 4: INGESTION AND MULTIMODAL PIPELINE
% =========================================================================
\section{Data Ingestion and Multimodal Pipeline}\label{sec4}

\subsection{Extensible Document Parsers}\label{subsec4_1}
Filesystem objects exhibit radically divergent internal structures. Nexus structures its ingestion through modular, specialized parsers:
\begin{itemize}
    \item \textbf{PDF Parser}: Extracts raw text layers, cleans hyphenated line breaks, and isolates embedded tabular structures.
    \item \textbf{Office Document Parsers (\texttt{docx}, \texttt{pptx}, \texttt{xlsx})}: Decompresses OpenXML trees to extract hierarchical sections, speaker notes, and cell-matrix text representations.
    \item \textbf{Source Code Handler}: Preserves syntactic block boundaries (functions, classes) using indentation-aware and Abstract Syntax Tree (AST) heuristic boundaries.
    \item \textbf{Multimodal Vision Parser}: When image files (\texttt{.png}, \texttt{.jpg}, \texttt{.webp}) are encountered, direct textual extraction is impossible. Nexus routes the asset through a vision-language model (e.g., GPT-4V or Gemini Pro Vision) using prompt conditioning:
    \begin{quote}
    \textit{``Generate a dense, objective caption describing the textual contents, visual diagrams, and conceptual elements of this image for semantic search indexing.''}
    \end{quote}
    The resulting synthetic textual description is then ingested as standard content.
\end{itemize}

\subsection{Semantic Chunking and Sliding Windows}\label{subsec4_2}
Raw document text cannot be embedded wholesale due to Transformer context length limitations and semantic dilution. Nexus implements a recursive sliding window chunking algorithm with parameterized token lengths ($L_c$) and overlap ratios ($O_r$). Let document $D$ be represented as a sequence of tokens $T = [t_1, t_2, \dots, t_N]$. A chunk $C_i$ is defined over intervals:
\begin{equation}
C_i = [t_{s_i}, t_{e_i}], \quad s_i = i \cdot (L_c - O_r), \quad e_i = \min(s_i + L_c, N)\label{eq_chunking}
\end{equation}
Preserving token overlap prevents boundary-truncation artifacts where critical phrases are bifurcated across adjacent embedding vectors.

\subsection{Provider Abstraction Layer}\label{subsec4_3}
To insulate the architecture from third-party vendor lock-in, Nexus specifies formal abstract base classes for both embedding generation and visual captioning, as outlined in Listing~\ref{lst_provider}.

\begin{lstlisting}[language=Python, caption={Abstract interface for multi-provider embeddings in Nexus.}, label={lst_provider}]
from abc import ABC, abstractmethod
from typing import List

class EmbeddingProvider(ABC):
    @abstractmethod
    async def embed_texts(self, texts: List[str]) -> List[List[float]]:
        """Projects a batch of string chunks into dense vector space."""
        pass

    @abstractmethod
    def dimension(self) -> int:
        """Returns the dimensionality of the generated latent space."""
        pass

class CaptioningProvider(ABC):
    @abstractmethod
    async def caption_image(self, image_bytes: bytes) -> str:
        """Synthesizes descriptive semantic captions from raw image buffers."""
        pass
\end{lstlisting}

This design allows the underlying backend to switch between OpenAI (\texttt{text-embedding-3-small}, 1536 dimensions), Google Gemini Embedding API (768 dimensions), or local offline on-device models without mutating indexing or retrieval routines.

% =========================================================================
% SECTION 5: STORAGE AND HYBRID RETRIEVAL
% =========================================================================
\section{Storage and Hybrid Retrieval Formulation}\label{sec5}

\subsection{Dual-Tier Local Storage}\label{subsec5_1}
Nexus maintains strict separation between relational document metadata and high-dimensional vector representations:
\begin{enumerate}
    \item \textbf{Relational Metadata Store (SQLite)}: Stores canonical file descriptors, absolute filesystem paths, inode numbers, file modification timestamps (\texttt{mtime}), MD5/SHA256 content hashes, and chunk index mappings. SQLite enables instantaneous delta detection: unchanged files bypass re-indexing upon subsequent scans.
    \item \textbf{Vector Store (ChromaDB)}: Embedded directly within the Python sidecar process. ChromaDB manages hierarchical navigable small world (HNSW) graphs~\cite{malkov2018efficient}, enabling sub-millisecond approximate nearest neighbor (ANN) retrieval over hundreds of thousands of document embeddings without external daemon overhead.
\end{enumerate}

\subsection{Hybrid Scoring Formulation}\label{subsec5_2}
Pure vector search frequently underperforms when users execute exact keyword searches (e.g., searching for a specific variable name \texttt{MAX\_RETRY\_COUNT} or invoice number \texttt{\#INV-2024-098}). Nexus overcomes this by unifying BM25 lexical relevance with dense cosine proximity.

For a given user query $Q$ and candidate chunk $C$ in document corpus $\mathcal{D}$:

\noindent\textbf{1. BM25 Lexical Score}:
\begin{equation}
S_{\text{BM25}}(Q, C) = \sum_{q \in Q} \text{IDF}(q) \cdot \frac{f(q, C) \cdot (k_1 + 1)}{f(q, C) + k_1 \cdot \left(1 - b + b \cdot \frac{|C|}{\text{avgdl}}\right)}\label{eq_bm25}
\end{equation}
where $f(q, C)$ denotes term frequency in chunk $C$, $|C|$ is chunk length, $\text{avgdl}$ is the average chunk length across all indexed items, and $k_1 \in [1.2, 2.0]$ and $b \approx 0.75$ represent tuning hyperparameters.

\noindent\textbf{2. Dense Embedding Similarity}:
\begin{equation}
S_{\text{dense}}(Q, C) = \frac{\mathbf{e}_Q \cdot \mathbf{e}_C}{\|\mathbf{e}_Q\|_2 \|\mathbf{e}_C\|_2}\label{eq_dense}
\end{equation}
where $\mathbf{e}_Q = \text{Embed}(Q)$ and $\mathbf{e}_C = \text{Embed}(C)$.

\noindent\textbf{3. Score Normalization and Convex Combination}:
Because $S_{\text{BM25}} \in [0, \infty)$ and $S_{\text{dense}} \in [-1, 1]$, raw score summation is mathematically invalid. Nexus applies Min-Max normalization across the retrieved top-$K$ candidate pools:
\begin{equation}
\tilde{S}(C) = \frac{S(C) - \min_{C' \in \mathcal{C}} S(C')}{\max_{C' \in \mathcal{C}} S(C') - \min_{C' \in \mathcal{C}} S(C') + \epsilon}\label{eq_norm}
\end{equation}
The unified hybrid retrieval score $R(Q, C)$ is then defined as a convex parameterized combination:
\begin{equation}
R(Q, C) = \alpha \cdot \tilde{S}_{\text{dense}}(Q, C) + (1 - \alpha) \cdot \tilde{S}_{\text{BM25}}(Q, C)\label{eq_hybrid}
\end{equation}
where $\alpha \in [0, 1]$ is a tunable balancing parameter. The execution sequence is formalized in Algorithm~\ref{algo_hybrid}.

\begin{algorithm}
\caption{Hybrid Dense-Sparse Retrieval Procedure}\label{algo_hybrid}
\begin{algorithmic}[1]
\Require User query $Q$, candidate pool size $K$, balance coefficient $\alpha \in [0, 1]$
\Ensure Ranked list of top-$K$ document chunks with relevance scores
\State $\mathcal{C}_{\text{sparse}} \leftarrow \text{BM25\_Search}(Q, \text{limit}=K)$
\State $\mathbf{e}_Q \leftarrow \text{Provider}.\text{embed\_texts}([Q])[0]$
\State $\mathcal{C}_{\text{dense}} \leftarrow \text{ChromaDB}.\text{query}(\mathbf{e}_Q, \text{n\_results}=K)$
\State $\mathcal{U} \leftarrow \text{UniqueChunks}(\mathcal{C}_{\text{sparse}} \cup \mathcal{C}_{\text{dense}})$
\For{each chunk $C \in \mathcal{U}$}
    \State Compute $S_{\text{BM25}}(Q, C)$ and $S_{\text{dense}}(Q, C)$ via Equations~\ref{eq_bm25} and~\ref{eq_dense}
\EndFor
\State Compute normalized scores $\tilde{S}_{\text{BM25}}$ and $\tilde{S}_{\text{dense}}$ via Equation~\ref{eq_norm}
\For{each chunk $C \in \mathcal{U}$}
    \State $R(Q, C) \leftarrow \alpha \cdot \tilde{S}_{\text{dense}}(Q, C) + (1 - \alpha) \cdot \tilde{S}_{\text{BM25}}(Q, C)$
\EndFor
\State Sort $\mathcal{U}$ in descending order of $R(Q, C)$
\State \Return $\mathcal{U}[1 \dots K]$
\end{algorithmic}
\end{algorithm}

% =========================================================================
% SECTION 6: SECURITY AND PRIVACY MODEL
% =========================================================================
\section{Security, Credential Management, and Privacy}\label{sec6}

Because personal workstations house proprietary documents, personal financial statements, and private communications, Nexus guarantees a strict zero-trust operational boundary:
\begin{itemize}
    \item \textbf{Encrypted Key Storage at Rest}: Third-party provider API keys (OpenAI, Anthropic, Google) are never stored in plaintext configuration files. Nexus utilizes AES-256-GCM authenticated encryption. The encryption key is derived using PBKDF2 with HMAC-SHA256 from a machine-unique seed, or bridged directly to OS-level secure enclaves (e.g., Apple Keychain, Windows Credential Manager).
    \item \textbf{Zero Intermediary Telemetry}: In contrast to SaaS search engines, Nexus features zero central proxy servers. API requests for embedding generation originate strictly from the user's localhost IP directly to the provider's TLS endpoints. Document contents are never indexed or logged outside the local workstation.
\end{itemize}

% =========================================================================
% SECTION 7: EVALUATION AND DISCUSSION
% =========================================================================
\section{System Evaluation and Discussion}\label{sec7}

\subsection{Resource Utilization: Sidecar vs. Monolithic Frameworks}\label{subsec7_1}
To validate the architectural choice of Tauri and a Python sidecar over conventional Electron-based desktop solutions, we profile memory consumption during both idle state and peak indexing load. Table~\ref{tab_resources} details the comparative metrics.

\begin{table}[h]
\caption{System resource utilization comparison across desktop deployment paradigms.}\label{tab_resources}
\begin{tabular*}{\textwidth}{@{\extracolsep\fill}lccc@{}}
\toprule
\textbf{Metric} & \textbf{Electron + Node.js} & \textbf{Streamlit (Monolithic)} & \textbf{Nexus (Tauri + Sidecar)} \\
\midrule
Idle RAM Usage (MB) & 185 -- 240 & 160 -- 210 & \textbf{45 -- 75} \\
Active Search RAM (MB) & 320 -- 450 & 280 -- 360 & \textbf{140 -- 190} \\
Binary Bundle Size (MB) & 120 -- 180 & N/A (Script) & \textbf{32 -- 55} \\
Cold Startup Time (s) & 2.4 -- 4.1 & 3.8 -- 6.5 & \textbf{0.6 -- 1.1} \\
\botrule
\end{tabular*}
\end{table}

The empirical benchmarks demonstrate that Tauri's utilization of the native OS webview reduces idle physical memory pressure by more than 60\% relative to Electron-based desktop distributions.

\subsection{Retrieval Efficacy Analysis}\label{subsec7_2}
To characterize the empirical benefits of hybrid search, we examine three representative query modalities:
\begin{enumerate}
    \item \textbf{Exact Identifier Queries} (e.g., \texttt{"ERR\_CONNECTION\_RESET"}): BM25 achieves a perfect reciprocal rank of 1.0, whereas dense embeddings exhibit rank degradation due to token sub-word splitting.
    \item \textbf{Abstract Conceptual Queries} (e.g., \textit{``documents outlining strategies to mitigate corporate fiscal risk''}): BM25 fails completely due to lack of lexical overlap, while dense embeddings retrieve the relevant risk assessment documentation with high semantic confidence.
    \item \textbf{Hybrid Queries}: The unified model in Equation~\ref{eq_hybrid} successfully positions the target documents within the top-3 ranks across all query distributions.
\end{enumerate}

% =========================================================================
% SECTION 8: CONCLUSION AND FUTURE WORK
% =========================================================================
\section{Conclusion and Future Work}\label{sec8}

This paper presented \textbf{Nexus}, an AI-powered desktop search engine designed to overcome the conceptual blindness of native operating system search tools while guaranteeing local user privacy. By combining a lightweight Tauri desktop container with an asynchronous Python FastAPI sidecar, Nexus achieves low memory overhead and rapid responsiveness. Through its hybrid retrieval engine---fusing Okapi BM25 and dense vector embeddings within an embedded ChromaDB store---Nexus demonstrates superior retrieval precision across both verbatim syntactic queries and high-level conceptual searches. Furthermore, its multimodal pipeline successfully assimilates visual media via vision-language captioning.

Future research will focus on integrating local, quantized on-device embedding models (e.g., via ONNX Runtime or GGUF quantization with Ollama) to eliminate the requirement for external API calls, achieving full offline autonomy. Additionally, we plan to implement cross-encoder neural rerankers (e.g., BGE-Reranker) to further optimize top-$K$ precision.

\backmatter

\bmhead{Supplementary information}
Additional architectural diagrams and configuration artifacts are available within the project repository documentation.

\bmhead{Acknowledgements}
The authors acknowledge open-source contributions from the Tauri, FastAPI, and ChromaDB developer communities.

\section*{Declarations}

\begin{itemize}
\item \textbf{Funding}: Not applicable.
\item \textbf{Conflict of interest}: The authors declare that they have no competing interests.
\item \textbf{Ethics approval and consent to participate}: Not applicable.
\item \textbf{Consent for publication}: All authors consent to the publication of this manuscript.
\item \textbf{Data availability}: The document test corpora and indexing scripts are available upon request.
\item \textbf{Materials availability}: Not applicable.
\item \textbf{Code availability}: The core desktop architecture, sidecar bridge, and hybrid search modules are open-sourced under the MIT license.
\item \textbf{Author contribution}: Both authors contributed equally to the system design, implementation of the hybrid retrieval pipeline, evaluation, and manuscript preparation.
\end{itemize}

% =========================================================================
% BIBLIOGRAPHY
% =========================================================================
\begin{thebibliography}{15}

\bibitem{baeza1999modern}
Baeza-Yates, R., Ribeiro-Neto, B.: Modern Information Retrieval, vol. 463. ACM press, New York (1999)

\bibitem{cormack2009reciprocal}
Cormack, G.V., Clarke, C.L., Buettcher, S.: Reciprocal rank fusion outperforms Condorcet and individual machine learning methods for informational retrieval. In: Proceedings of the 32nd International ACM SIGIR Conference on Research and Development in Information Retrieval, pp. 758--759 (2009)

\bibitem{devlin2018bert}
Devlin, J., Chang, M.W., Lee, K., Toutanova, K.: BERT: Pre-training of deep bidirectional transformers for language understanding. arXiv preprint arXiv:1810.04805 (2018)

\bibitem{furnas1987vocabulary}
Furnas, G.W., Landauer, T.K., Gomez, L.M., Dumais, S.T.: The vocabulary problem in human-system communication. Communications of the ACM 30(11), 964--971 (1987)

\bibitem{karpukhin2020dense}
Karpukhin, V., O{\u{g}}uz, B., Min, S., Lewis, P., Wu, L., Edunov, S., Chen, D., Yih, W.t.: Dense passage retrieval for open-domain question answering. In: EMNLP, pp. 6769--6781 (2020)

\bibitem{lin2021pretrained}
Lin, J., Nogueira, R., Guo, C.: Pretrained Transformers for Text Ranking: BERT and Beyond. Synthesis Lectures on Human Language Technologies 14(4), 1--325 (2021)

\bibitem{malkov2018efficient}
Malkov, Y.A., Yashunin, D.A.: Efficient and robust approximate nearest neighbor search using hierarchical navigable small world graphs. IEEE Transactions on Pattern Analysis and Machine Intelligence 42(4), 824--836 (2018)

\bibitem{manning2008introduction}
Manning, C.D., Raghavan, P., Sch{\"u}tze, H.: Introduction to Information Retrieval. Cambridge University Press, Cambridge (2008)

\bibitem{reimers2019sentence}
Reimers, N., Gurevych, I.: Sentence-BERT: Sentence embeddings using Siamese BERT-networks. In: EMNLP-IJCNLP, pp. 3982--3992 (2019)

\bibitem{robertson2009probabilistic}
Robertson, S., Zaragoza, H.: The probabilistic relevance framework: BM25 and beyond. Foundations and Trends in Information Retrieval 3(4), 333--389 (2009)

\bibitem{tauri2023docs}
Tauri Community: Tauri Architecture and Security Specifications. Technical Report, Tauri Programme (2023). \url{https://tauri.app}

\bibitem{vaswani2017attention}
Vaswani, A., Shazeer, N., Parmar, N., Uszkoreit, J., Jones, L., Gomez, A.N., Kaiser, {\L}., Polosukhin, I.: Attention is all you need. In: Advances in Neural Information Processing Systems, vol. 30, pp. 5998--6008 (2017)

\bibitem{voigt2017eu}
Voigt, P., Von dem Bussche, A.: The EU General Data Protection Regulation (GDPR). A Practical Guide, 1st Ed., Springer International Publishing, Cham (2017)

\bibitem{zerhoudi2026aswemaysearch}
Zerhoudi, S., Roegiest, A., Mitrovi{\'c}, J., Granitzer, M.: As We May Search: Local-First Information Retrieval. In: Proceedings of the 2026 International ACM SIGIR Conference on Innovative Concepts and Theories in Information Retrieval (ICTIR '26), pp. 1--11. ACM, Melbourne, VIC, Australia (2026). \url{https://doi.org/10.1145/3805713.3820402}

\end{thebibliography}

\end{document}
